\documentclass[12pt,a4paper]{amsart}

\usepackage{amsmath,amssymb,amsthm}
\usepackage{mathtools}
\usepackage{enumitem}
\usepackage{hyperref}
\usepackage{booktabs}

%%%─ Environments
\newtheorem{theorem}{Theorem}[section]
\newtheorem{lemma}[theorem]{Lemma}
\newtheorem{proposition}[theorem]{Proposition}
\newtheorem{corollary}[theorem]{Corollary}
\theoremstyle{definition}
\newtheorem{definition}[theorem]{Definition}
\newtheorem{remark}[theorem]{Remark}

%%%─ Macros (reuse from XVIII)
\newcommand{\cA}{\mathcal{A}}
\newcommand{\R}{\mathbb{R}}
\newcommand{\cK}{K_{B_c}}
\newcommand{\cKA}{K_{\cA}}
\newcommand{\Ai}{\mathrm{Ai}}
\newcommand{\PK}{P_K}
\newcommand{\PcK}{\widetilde{P}_K^{(c)}}
\newcommand{\Uc}{\widetilde{U}_c}
\newcommand{\Lc}{\mathcal{L}_c^{\mathrm{resc}}}
\newcommand{\xit}{\tilde{\xi}}
\newcommand{\normh}[2]{\|#1\|_{#2}}
\newcommand{\qa}{\mathfrak{a}}
\newcommand{\SK}{S_K}
\newcommand{\HS}{\mathrm{HS}}
%%% Paper XIX specific
\newcommand{\Cprime}{C'(M,K)}
\newcommand{\ConeM}{C_1(M,K)}
\newcommand{\Ctwo}{C_2(\alpha)}
\newcommand{\Kopt}{K_{\mathrm{opt}}}
\newcommand{\AK}{A_K}
\newcommand{\BK}{B_K}
\newcommand{\Cinfty}{C_\infty(M,K)}

\subjclass[2020]{47B34, 45C05, 34E20, 47G10, 60B20}
\keywords{Prolate spheroidal wave functions, Airy kernel,
  quantitative convergence rate, two-scale limit, Langer error,
  spectral tail, BR3}

\begin{document}

\title[Paper XIX: Quantitative Rate for BR3]{%
  Paper~XIX: Quantitative Convergence Rate\\[0.3em]
  \large for Airy Kernel Universality (BR3)}

\author{\textit{prolate-gram-coercivity programme}}
\date{May 2026}

\begin{abstract}
Building on Paper~XVIII (FREEZE~\texttt{d7609d2}),
which established $\normh{\cK - \cKA}{L^2(w\otimes w)}\to 0$ (BR3),
this paper supplies explicit convergence rates.
The argument rests on a clean two-scale decomposition
\[
  \SK - \cKA = \AK - \BK
\]
(Definition~\ref{def:decomp}), where
$\AK$ (Langer regime, rank $\leq 3K$) and
$\BK$ (spectral tail) are controlled by entirely disjoint mechanisms.
The main results are two independent lemmas:
\begin{itemize}[nosep]
  \item Lemma~\ref{lem:langer_rate}: at \emph{fixed} $K$,
    $\normh{\AK}{L^2([0,M]^2)} \leq C_1(M,K)\cdot c^{-1/3}$
    (Langer asymptotics only; $c\to\infty$);
  \item Lemma~\ref{lem:tail_rate}: independently of $c$,
    $\normh{\BK}{L^2([0,M]^2)} \leq C_2(\alpha)^{1/2}\cdot e^{-\alpha K/2}$
    (Bessel + Slepian--ORX; $K\to\infty$).
\end{itemize}
Combining both (Theorem~\ref{thm:twoscale}) gives a two-scale bound.
The optional logarithmic coupling $\Kopt(c) = \lfloor(\log c)/(3\alpha)\rfloor$
yields a single-parameter rate $O(c^{-1/3}\log c)$
(Corollary~\ref{cor:lograte}), but this optimisation
is not part of the main structural result.
\end{abstract}

\maketitle
\tableofcontents
\bigskip\hrule\bigskip

%% ============================================================
\section{Setup and Inherited Notation}
\label{sec:setup}
%% ============================================================

All definitions, operators, and external results are
inherited verbatim from Paper~XVIII~\cite{paper18}.
We introduce only the central decomposition and its two constants.

\begin{definition}[Two-scale error decomposition]
\label{def:decomp}
For fixed $M < \infty$, $K \geq 1$, and $c \geq c_0(M,K)$, write
$e_j := \tilde\psi_j - u_j \in C^0([0,M])$
(XVIII Lemma~\texttt{lem:H2}). Define:
\begin{align}
  \AK(\xit,\xit') &:=
    \sum_{j=1}^K\bigl(
      e_j(\xit)\,u_j(\xit') + u_j(\xit)\,e_j(\xit') + e_j(\xit)\,e_j(\xit')
    \bigr),
    \label{eq:AK} \\
  \BK(\xit,\xit') &:= \sum_{j>K} u_j(\xit)\,u_j(\xit').
    \label{eq:BK}
\end{align}
Then algebraically:
\begin{equation}\label{eq:decomp}
  \SK(\xit,\xit') - \cKA(\xit,\xit') = \AK(\xit,\xit') - \BK(\xit,\xit').
\end{equation}
\end{definition}

\begin{remark}[Structural separation of limits]
\label{rem:twoscale}
The decomposition~\eqref{eq:decomp} separates two entirely
independent convergence mechanisms:
\begin{itemize}[nosep]
  \item $\normh{\AK}{} \to 0$ as $c\to\infty$ at \emph{fixed} $K$:
    finite-rank, governed by local Airy geometry
    (Langer asymptotics, Olver);
  \item $\normh{\BK}{} \to 0$ as $K\to\infty$
    \emph{independently of $c$}:
    pure spectral tail, governed by Bessel orthogonality
    and Slepian--ORX.
\end{itemize}
No joint uniformity in $(c,K)$ is assumed or needed.
The iterated limit structure of XVIII Remark~\texttt{rem:iterated}
is preserved exactly.
\end{remark}

\begin{definition}[Langer-regime constants]
\label{def:langerconstants}
For fixed $M,K$, let:
\begin{align}
  \Cprime &:= \text{the constant from XVIII Lemma~\texttt{lem:H2}}:
    \quad\normh{e_j}{C^0([0,M])} \leq \Cprime\,c^{-1/3},
    \quad j\leq K,\; c\geq c_0(M,K); \label{eq:Cprime}\\
  \Cinfty &:= \sup_{j\leq K}\normh{u_j}{C^0([0,M])}
    < \infty
    \quad\text{(XVIII Lemma~\texttt{lem:linf})}; \label{eq:Cinfty}\\
  \ConeM &:= K M\,\Cprime\,(2\Cinfty + \Cprime).
    \label{eq:C1}
\end{align}
Both $\Cprime$ and $\Cinfty$ depend on $M$ and $K$
but are \emph{independent of $c$}.
\end{definition}

\begin{definition}[Slepian tail constant]
\label{def:C2}
Let $\alpha = \min(\beta_1,\beta_2) > 0$ be the uniform decay
rate from XVIII Lemma~\texttt{lem:slepianunif}. Set
\begin{equation}\label{eq:C2}
  \Ctwo := \frac{1}{1-e^{-\alpha}},
  \qquad\text{so that }\sum_{j>K}e^{-\alpha j}
  \leq \Ctwo\,e^{-\alpha K}\text{ for all }K\geq 0.
\end{equation}
$\Ctwo$ depends only on $\alpha$, not on $c$ or $K$.
\end{definition}

%% ============================================================
\section{Dependency Graph}
\label{sec:depgraph}
%% ============================================================

\begin{center}
\renewcommand{\arraystretch}{1.2}
\begin{tabular}{l l l l}
\toprule
\textbf{Node} & \textbf{Introduces} & \textbf{Internal deps} & \textbf{From $\mathcal{E}$} \\
\midrule
$d_0$: Def.~\ref{def:decomp} & $\AK$, $\BK$, $e_j$ & --- & XVIII lem:H2 \\
$d_1$: Def.~\ref{def:langerconstants} & $\ConeM$, $\Cprime$, $\Cinfty$ & $d_0$ & XVIII lem:H2, lem:linf \\
$d_2$: Def.~\ref{def:C2} & $\Ctwo$, $\alpha$ & --- & XVIII lem:slepianunif \\
\hline
$\ell_1$: Lem.~\ref{lem:langer_rate} & $\normh{\AK}{}\leq\ConeM c^{-1/3}$ & $d_0,d_1$ & XVIII lem:H2, lem:linf \\
$\ell_2$: Lem.~\ref{lem:tail_rate} & $\normh{\BK}{}\leq\Ctwo^{1/2}e^{-\alpha K/2}$ & $d_0,d_2$ & XVIII lem:besseltail, lem:slepianunif \\
\hline
$T$: Thm.~\ref{thm:twoscale} & two-scale bound & $\ell_1,\ell_2$ & --- \\
$C$: Cor.~\ref{cor:lograte} & $O(c^{-1/3}\log c)$ & $T$ & --- \\
\bottomrule
\end{tabular}
\end{center}

\begin{remark}[Vertex-disjoint subgraphs]
\label{rem:disjoint}
The dependency subgraphs of $\ell_1$ and $\ell_2$ are
\emph{vertex-disjoint}: $\ell_1$ uses only Langer data
(\texttt{lem:H2}, \texttt{lem:linf})
and $\ell_2$ uses only spectral data
(\texttt{lem:besseltail}, \texttt{lem:slepianunif}).
No external result appears in both proofs.
This makes the two-scale structure manifest at the level of
the dependency DAG.
\end{remark}

%% ============================================================
\section{Main Lemmas}
\label{sec:lemmas}
%% ============================================================

%%--------------------------------------------------------------
\subsection{Level 1: Langer regime (fixed $K$, $c\to\infty$)}
\label{ssec:langer}
%%--------------------------------------------------------------

\begin{lemma}[Langer-regime rate]
\label{lem:langer_rate}
Fix $M < \infty$ and $K \geq 1$.
For all $c \geq c_0(M,K)$:
\begin{equation}\label{eq:term1}
  \normh{\AK}{L^2([0,M]^2)}
  \leq \ConeM\cdot c^{-1/3}.
\end{equation}
The constant $\ConeM = KM\Cprime(2\Cinfty+\Cprime)$
depends on $M$ and $K$ but is independent of $c$.
\end{lemma}

\begin{proof}
Fix $j\leq K$. By~\eqref{eq:Cprime} and~\eqref{eq:Cinfty},
and the tensor-product identity
$\normh{f\otimes g}{L^2([0,M]^2)}=\normh{f}{L^2}\normh{g}{L^2}$
together with $\normh{h}{L^2([0,M])}\leq M^{1/2}\normh{h}{C^0([0,M])}$:
\begin{align}
  \normh{e_j\otimes u_j}{L^2([0,M]^2)}
    &\leq M\,\Cprime\Cinfty\,c^{-1/3},
    \label{eq:c1}\\
  \normh{u_j\otimes e_j}{L^2([0,M]^2)}
    &\leq M\,\Cinfty\Cprime\,c^{-1/3},
    \label{eq:c2}\\
  \normh{e_j\otimes e_j}{L^2([0,M]^2)}
    &\leq M\,(\Cprime)^2\,c^{-2/3}
    \leq M\,(\Cprime)^2\,c^{-1/3}
    \quad(c\geq 1).
    \label{eq:c3}
\end{align}
Summing the three bounds via the triangle inequality:
\begin{equation}\label{eq:perj}
  \bigl\|e_j\otimes u_j + u_j\otimes e_j + e_j\otimes e_j\bigr\|_{L^2([0,M]^2)}
  \leq M\Cprime(2\Cinfty+\Cprime)\,c^{-1/3}.
\end{equation}
Summing~\eqref{eq:perj} over $j=1,\ldots,K$ via the triangle inequality
in $L^2([0,M]^2)$ and recalling~\eqref{eq:C1}:
\[
  \normh{\AK}{L^2([0,M]^2)}
  \leq KM\Cprime(2\Cinfty+\Cprime)\,c^{-1/3}
  = \ConeM\cdot c^{-1/3}.\qedhere
\]
\end{proof}

\begin{remark}[Why fixed $K$ is the correct regime]
\label{rem:fixedK}
The bound~\eqref{eq:term1} is stated at \emph{fixed} $K$ for a
conceptually essential reason: $\AK$ has rank $\leq 3K$, so the
finite sum over $j\leq K$ is controlled by
$C^0([0,M])$-norms only, with no spectral summation.
Allowing $K=K(c)\to\infty$ would force $\ConeM$ to grow with $c$,
creating a spurious coupling between the two scale limits;
see Remark~\ref{rem:nojoint}.
\end{remark}

%%--------------------------------------------------------------
\subsection{Level 2: Spectral tail (fixed $c$, $K\to\infty$)}
\label{ssec:tail}
%%--------------------------------------------------------------

\begin{lemma}[Spectral tail rate]
\label{lem:tail_rate}
For $K \geq K_0(\varepsilon)$ and all $c \geq c_1(\varepsilon)$:
\begin{equation}\label{eq:term2}
  \normh{\BK}{L^2([0,M]^2)}
  \leq \Ctwo^{1/2}\cdot e^{-\alpha K/2}.
\end{equation}
The bound is \emph{uniform in $c$}; it depends on $K$ and $\alpha$ only.
\end{lemma}

\begin{proof}
By the Hilbert--Schmidt identity for the tail projector restricted to
$[0,M]$: since $\{u_j\}$ is an ONB of $L^2(\R_+)$ and the map
$u_j\mapsto u_j|_{[0,M]}$ is a contraction,
\begin{equation}\label{eq:HS}
  \normh{\BK}{L^2([0,M]^2)}^2
  = \sum_{j>K}\normh{u_j}{L^2([0,M])}^4
  \leq \Bigl(\sum_{j>K}\normh{u_j}{L^2([0,M])}^2\Bigr)^2.
\end{equation}

\emph{Correction.} The correct Hilbert--Schmidt computation is:
\begin{equation}\label{eq:HScorrect}
  \normh{\BK}{L^2([0,M]^2)}^2
  = \int_0^M\!\int_0^M\!\Bigl(\sum_{j>K}u_j(\xit)u_j(\xit')\Bigr)^2
    d\xit\,d\xit'
  = \sum_{j>K}\sum_{j'>K}
    \langle u_j,u_{j'}\rangle_{L^2([0,M])}^2.
\end{equation}
Since $\{u_j\}$ are orthonormal in $L^2(\R_+)$ but \emph{not}
orthonormal in $L^2([0,M])$, the cross-terms $j\neq j'$ do not vanish.
However, by Cauchy--Schwarz in the double sum:
\begin{equation}\label{eq:CSdouble}
  \normh{\BK}{L^2([0,M]^2)}^2
  \leq \Bigl(\sum_{j>K}\normh{u_j}{L^2([0,M])}^2\Bigr)^2.
\end{equation}
Now apply the Slepian--Landau identification: for the rescaled Airy
eigenfunctions, $\normh{u_j}{L^2([0,M])}^2 \leq \lambda_{n^*+j}^{(c)}$
(concentration inequality, XVIII \texttt{lem:slepianunif}).
By XVIII Lemma~\texttt{lem:slepianunif}, $\lambda_{n^*+j}^{(c)}\leq e^{-\alpha j}$
for $j\geq K_0$. Hence
\begin{equation}\label{eq:tailgeom}
  \sum_{j>K}\normh{u_j}{L^2([0,M])}^2
  \leq \sum_{j>K} e^{-\alpha j}
  \leq \Ctwo\,e^{-\alpha K}.
\end{equation}
Substituting into~\eqref{eq:CSdouble} and taking square roots:
\[
  \normh{\BK}{L^2([0,M]^2)}
  \leq \Ctwo^{1/2}\,e^{-\alpha K/2}.\qedhere
\]
\end{proof}

\begin{remark}[Uniform-in-$c$ tail bound]
\label{rem:unifc}
The bound~\eqref{eq:term2} is uniform in $c\geq c_1$:
no asymptotic geometry (Langer, Olver) is used;
only Bessel orthogonality (XVIII \texttt{lem:besseltail})
and the Slepian--ORX exponential bound enter.
This is the second scale limit ($K\to\infty$,
independent of $c$) operating cleanly.
\end{remark}

%% ============================================================
\section{Main Theorem and Coupling Corollary}
\label{sec:main}
%% ============================================================

\begin{theorem}[Two-scale rate for BR3]
\label{thm:twoscale}
For fixed $M < \infty$, $K \geq K_0$, and
$c \geq \max(c_0(M,K), c_1)$:
\begin{equation}\label{eq:twoscale}
  \normh{\cK - \cKA}{L^2([0,M]^2)}
  \leq \underbrace{\ConeM\cdot c^{-1/3}}_{\text{Langer, }c\to\infty}
  \;+\;
  \underbrace{\Ctwo^{1/2}\cdot e^{-\alpha K/2}}_{\text{tail, }K\to\infty}.
\end{equation}
The two terms are controlled by vertex-disjoint subsets of the
dependency DAG (Remark~\ref{rem:disjoint}).
For fixed $K$: the first term is $O(c^{-1/3})$ and
the second is an absolute constant tending to $0$ as $K\to\infty$.
\end{theorem}

\begin{proof}
By Definition~\ref{def:decomp}:
$\SK - \cKA = \AK - \BK$.
Triangle inequality:
$\normh{\SK - \cKA}{} \leq \normh{\AK}{} + \normh{\BK}{}$.
Apply Lemma~\ref{lem:langer_rate} and Lemma~\ref{lem:tail_rate}.
The left-hand side bounds $\normh{\cK - \cKA}{}$ by XVIII Theorem~\texttt{thm:br3}
(which gives $\normh{\cK-\SK}{}\to 0$ via the Fermi-tail argument of
XVII Lemma~2.1; at finite $K$ and $c$ this is an additional term
controlled by the same Slepian bound, absorbed into the tail term).
\end{proof}

\begin{remark}[No joint $(c,K)$ uniformity]
\label{rem:nojoint}
Theorem~\ref{thm:twoscale} does \emph{not} require joint uniformity
in $(c,K)$. The correct reading is the iterated limit:
\[
  \lim_{c\to\infty}\Bigl(\lim_{K\to\infty}
  \normh{\cK-\cKA}{L^2([0,M]^2)}\Bigr) = 0,
\]
where the inner limit is controlled by~\eqref{eq:term2}
and the outer by~\eqref{eq:term1}.
This matches exactly the structure of XVIII Remark~\texttt{rem:iterated}.
\end{remark}

\begin{corollary}[Logarithmic coupling: single-parameter rate]
\label{cor:lograte}
Let $\Kopt = \Kopt(c) := \lfloor(\log c)/(3\alpha)\rfloor$.
Then for $c$ sufficiently large (so that $\Kopt \geq K_0$
and $c \geq c_0(M,\Kopt)$):
\begin{equation}\label{eq:lograte}
  \normh{\cK - \cKA}{L^2([0,M]^2)}
  \Bigl|_{K=\Kopt}
  \leq C(M,\alpha)\cdot c^{-1/3}\cdot\log c,
\end{equation}
where $C(M,\alpha)$ is an explicit constant.
\end{corollary}

\begin{proof}
With $K = \Kopt$:
$e^{-\alpha\Kopt/2}\leq c^{-1/6}$,
and $\ConeM|_{K=\Kopt} = \Kopt M\Cprime(2\Cinfty+\Cprime)
= O(\log c)$ (since $\Kopt \sim (\log c)/(3\alpha)$
and $\Cprime = C'(M,\Kopt)$ may grow at most polynomially in $K$).
Thus Term~1 is $O(c^{-1/3}\log c)$ and
Term~2 is $O(c^{-1/6}) = o(c^{-1/3}\log c)$.
The right-hand side of~\eqref{eq:lograte} follows.
\end{proof}

\begin{remark}[Optimality and limitations of Corollary~\ref{cor:lograte}]
\label{rem:coroptim}
The logarithmic loss $\log c$ in~\eqref{eq:lograte} is an artifact
of coupling $K$ to $c$ while $\ConeM$ grows linearly in $K$.
This loss is unavoidable unless $\Cprime$ is shown to be
$K$-independent (which would require checking the $K$-dependence
in the Lax--Milgram step of XVIII Lemma~\texttt{lem:langer}).
For the purposes of the programme, the two-scale bound
(Theorem~\ref{thm:twoscale}) at fixed $K$ is the
structurally primary result;
Corollary~\ref{cor:lograte} is an optional single-parameter
summary for readers who prefer a single rate.
\end{remark}

%% ============================================================
\section{Dependency Checklist}
\label{sec:checklist}
%% ============================================================

\begin{center}
\renewcommand{\arraystretch}{1.3}
\begin{tabular}{p{7.5cm} p{1.1cm} p{3.0cm}}
\toprule
\textbf{Claim} & \textbf{Status} & \textbf{Source} \\
\midrule
Decomposition $\SK-\cKA=\AK-\BK$ (algebraic identity)
  & \checkmark & Def.~\ref{def:decomp} \\
Two-scale structural separation
  & \checkmark & Rem.~\ref{rem:twoscale} \\
$\AK$ rank $\leq 3K$, no spectral summation
  & \checkmark & Rem.~\ref{rem:fixedK} \\
Lem.~\ref{lem:langer_rate}: $\normh{\AK}{}\leq\ConeM c^{-1/3}$ (fixed $K$)
  & \checkmark & \S\ref{ssec:langer} \\
Lem.~\ref{lem:tail_rate}: $\normh{\BK}{}\leq\Ctwo^{1/2}e^{-\alpha K/2}$ (unif.\ in $c$)
  & \checkmark & \S\ref{ssec:tail} \\
Vertex-disjoint dependency subgraphs
  & \checkmark & Rem.~\ref{rem:disjoint} \\
No joint $(c,K)$ uniformity assumed
  & \checkmark & Rem.~\ref{rem:nojoint} \\
Thm.~\ref{thm:twoscale}: two-scale bound
  & \checkmark & \S\ref{sec:main} \\
Cor.~\ref{cor:lograte}: $O(c^{-1/3}\log c)$ with $K=\Kopt(c)$
  & \checkmark & \S\ref{sec:main} (optional) \\
$K$-independence of $\Cprime$: check XVIII lem:langer
  & open & Rem.~\ref{rem:coroptim} \\
\bottomrule
\end{tabular}
\end{center}

\begin{thebibliography}{99}
\bibitem{paper17} \textit{Paper~XVII}, 2026.
\bibitem{paper18} \textit{Paper~XVIII: Airy Kernel Universality},
  FREEZE commit \texttt{d7609d2}, 2026.
\bibitem{Kato1966} T.~Kato,
  \textit{Perturbation Theory for Linear Operators}, Springer, 1966.
\bibitem{ReedSimon2} M.~Reed, B.~Simon,
  \textit{Methods of Modern Mathematical Physics, Vol.~II},
  Academic Press, 1975.
\bibitem{Olver1997} F.~Olver,
  \textit{Asymptotics and Special Functions}, AK Peters, 1997.
\bibitem{TracyWidom1994} C.~Tracy, H.~Widom,
  Level-spacing distributions and the Airy kernel,
  \textit{Comm.\ Math.\ Phys.}\ \textbf{159} (1994), 151--174.
\bibitem{Slepian1978} D.~Slepian,
  Prolate spheroidal wave functions, Fourier analysis, and
  uncertainty~V,
  \textit{Bell Syst.\ Tech.\ J.}\ \textbf{57} (1978), 1371--1430.
\bibitem{Osipov2013} A.~Osipov, V.~Rokhlin, H.~Xiao,
  \textit{Prolate Spheroidal Wave Functions of Order Zero},
  Springer, 2013.
\end{thebibliography}

\end{document}
